\section{Analysis and Calculations}

\subsection{Parameter Summary}
\begin{table}[H]
    \centering
    \small
    \begin{tabular}{l l l}
        \toprule
        \textbf{Parameter} & \textbf{Value} & \textbf{Comment} \\
        \midrule
        Nominal Input & \SI{12}{\volt} & Replace with the real operating point. \\
        Nominal Load Current & \SI{1.2}{\ampere} & Use measured or derived value. \\
        Switching Frequency & \SI{400}{\kilo\hertz} & Include if the design is switching-based. \\
        Thermal Ambient & \SI{50}{\celsius} & State the actual design assumption. \\
        \bottomrule
    \end{tabular}
    \caption{Example design parameters}
    \label{tab:parameters}
\end{table}

\subsection{Illustrative Sizing}
Use this section to document calculations that justify component values or
operating margins:

\begin{align}
    P_{out} &= V_{out} \times I_{out} \\
    \eta &= \frac{P_{out}}{P_{in}} \\
    P_{loss} &= P_{in} - P_{out}
\end{align}

\highlightbox{
Keep calculations reproducible. If a value comes from simulation, bench data, or
datasheet assumptions, say so explicitly.
}

\subsection{Design Notes}
\begin{itemize}
    \item explain why this topology or component family was selected;
    \item identify the limiting component or dominant loss path;
    \item summarize important design margins and unresolved assumptions.
\end{itemize}

\begin{figure}[H]
    \centering
    \begin{tikzpicture}
        \draw[brandgreen, line width=1pt] (0,0) rectangle (10,4.8);
        \draw[brandgreen, line width=1pt] (1.2,1.2) rectangle (3.8,3.6);
        \draw[brandgreen, line width=1pt] (6.2,1.2) rectangle (8.8,3.6);
        \draw[brandgreen, -latex, line width=1pt] (3.8,2.4) -- (6.2,2.4);
        \node at (2.5,2.4) {Source};
        \node at (7.5,2.4) {Load};
        \node[text width=8cm, align=center] at (5,0.5) {Placeholder for simplified power-path, control, or measurement flow diagram};
    \end{tikzpicture}
    \caption{Example schematic abstraction or analysis flow}
    \label{fig:analysis-placeholder}
\end{figure}
